% example.tex - an example workterm report that can be used as a template for writing NE 
% workterm reports

% We begin by calling the workreport class which includes all the definitions for the macros 
% and environments we will be using throughout the document.
\documentclass{workreport}

% We will use some packages to add some functionality to our script.
% The lipsum package is only used to create placeholder text for this example. It can be 
% deleted when your own text has been filled into the sections.
\usepackage{lipsum}
% The amsmath package provides various features to facilitate writing math formulas.
\usepackage{amsmath}

%%%%%%%%%%%%%%%%%%
%  IMPORTANT  INFORMATION  %
%%%%%%%%%%%%%%%%%%

% Fill in the title of your report 
\title{
	Example Work Term Report Style
}

% Fill in your name
\author{Regina Phalangey}

% Fill in your student number
\studentnumber{40440404}

% Fill in your student ID
\studentid{example}

% Fill in your previous academic term
\previousterm{2B}

% Fill in the name of your university
\university{University of Waterloo}

% Fill in your faculty
\faculty{Engineering}

% Fill in your program
\program{Nanotechnology Engineering}

% Fill in your program director
\programdirector{John Doe}

% Fill in the address of your university
\universityaddress{
	200 University Ave. West \\
	Waterloo, ON \\
	N2L 3G1
}

% Fill in your employers name
\employername{Employer Inc.}

% Fill in your employers address
\employeraddress{
	200 Payment St. \\
	Waterloo, ON \\
	A1A 1A1
}

% We will now begin writing the document
\begin{document}

%%%%%%%%%%%
% FRONT MATTER %
%%%%%%%%%%%

% The \begin{frontmatter} command will allow the use of the following environments:
% letter_of_submittal, contributions, summary, conclusions, recommendations, toctree
% (table of contents) and listoffiguresandtables (list of figures and list of tables).
\begin{frontmatter}

% The letter_of_submittal environment requires two inputs, the first being the greeting 
% and the second being the closing.
\begin{letter_of_submittal}{Dear Dr. Doe:}{Sincerely,}
    	
	% You're letter body.
	\lipsum[1-2]

\end{letter_of_submittal}

% The contribution environment only requires the contributions text.
\begin{contributions}

    	\lipsum[1-3]
	
\end{contributions}

% The summary environment only requires the summary text.
\begin{summary}

    	\lipsum[1-4]

\end{summary}

% The conclusions environment only requires the conclusions text.
\begin{conclusions}

        \lipsum[1-3]
        
\end{conclusions}

% The recommendations environment only requires the recmmendations text.
\begin{recommendations}

        \lipsum[1-3]
        
\end{recommendations}

% The following command generates a table of contents (toc) that automatically keeps 
% track of sections and their page numbers.
\toctree

% The following command generates a list of figures and a list of tables that
% automatically keep track of properly inserted figures and tables. The commands for 
% properly inserting figures and tables will be explained in the body section of this 
% document.
\listoffiguresandtables

\end{frontmatter}

%%%%%%
% BODY  %
%%%%%%

% The \begin{body} command will create the body environment. This environment 
% includes the formatting required for the body
\begin{body}

% The \section{title} command creates a section in the body of your report. You must
% specify the title of the section as the commands input. These sections headers will
% automatically be added to the table of contents.
\section{Introduction}

	\lipsum[1-6]

\section{Background}

	\lipsum[1-2]-

% The \subsection{title} command creates a subsection for the above mentioned
% section. You must specify the title of the section as the commands input. These 
% subsection headers will automatically be added to the table of contents, below the 
% proper section.
\subsection{A Subsection}

	\lipsum[1]

\subsection{A Second Subsection}

	\lipsum[2]

% The \subsubsection{title} command works exactly like the \section{title} and 
% \subsection{title} commands above. 
\subsubsection{A Subsection in a Subsection}

	\lipsum[1]

\subsubsection{Sub-ception}

	\lipsum[1]

	BWOOOONG!

\section{Data}

	\lipsum[1-2]
	
% The \ref{reference} command will create a reference to any \label command.
Here's a reference to table \ref{tbl:exampletable}

% The table environment is used to create a table. Note that by calling this environment, 
% this table will automatically be added to the list of tables.
\begin{table}
	
	% The \centering command centres the table.
	\centering
	
	% The tabular environment creates the table itself. For more information
	% and options when using this environment, refer to the link under Tables provided
	% in README.md.
	\begin{tabular}{|c|c|c|} \hline
		title & number & decimal \\ \hline
		heading & 1 & 1.7 \\ \hline
		another heading & 2 & 3.4 \\ \hline
	\end{tabular}
	% The \caption{title} command adds a caption to your table. This caption will then 
	% be used in the list of tables.
	\caption{A table demonstrating \LaTeX \, formatting}
	
	% The \label{title} command creates a reference that can be used to refer to this
	% table in the document. An example of this can be found on line 174.
	\label{tbl:exampletable}
	
\end{table}

\section{Math}

	\lipsum[1]

\subsection{The Bloch Equations}

	% The align environment is an important environment when writing out multiple 
	% mathematical equations. It is used when vertical alignment is desired. Usually 
	% binary relations, such as equal signs, are aligned.
	\begin{align}
		% The following are mathematical equations written using the amsmath 
		% package. For more information on writing mathematical equations in 
		% LaTeX, refer to the link under AMSMath in README.md.
		\frac{d\, M_x(t)}{dt} &= \gamma(\mathbf{M}(t) \times \mathbf{B}(t))_x - \frac{M_x(t)}{T_2} \\
		\frac{d\, M_y(t)}{dt} &= \gamma(\mathbf{M}(t) \times \mathbf{B}(t))_y - \frac{M_y(t)}{T_2} \\
		\frac{d\, M_z(t)}{dt} &= \gamma(\mathbf{M}(t) \times \mathbf{B}(t))_z - \frac{M_z(t) - M_0}{T_1}
	\end{align}

\subsection{The Schrodinger Equation}

	% The equation environment is used when only one equation needs to be written.
	\begin{equation}
		i\hbar \frac{\partial \, \Psi}{\partial t} = \hat{H}\Psi
	\end{equation}

\section{Graphics}

	Here's a stock image of a computer. It comes with a CC0 license, enabling
	free distribution for commercial and personal use. It's a JPEG.
	
	% The figure environment allows for the insertion of a figure into a LaTeX 
	% document. By creating this environment, the figure will be automatically be added
	% to the list of figures. For more information on inserting figures, refer to the link 
	% under Figures in README.md.
	\begin{figure}[!ht]
		\centering
		\label{fig:stock_computer}
		\includegraphics[width=5cm]{./stock-image.jpg}
		\caption{A stock computer, I just wanted a picture here}
	\end{figure}

\section{Conclusions}
	
	% The \gls{title} command allows the LaTeX document to keep track of where words
	% in the glossary are used. For more information on the glossary and acronyms, 
	% refer to the link under Glossary and Acronyms in README.md.
	\gls{LaTeX} is awesome.

\section{Recommendations}

	Use more \LaTeX. Also, use more \gls{Unix}. Also, citations like
    	\cite{schweiger_principles_2001} will lead to a better quality of work.

\end{body}

% The \newglossaryentry{title} command creates an entry for the glossary. It requires a 
% name and description. The title is used for making references to the glossary entry 
% throughout the document. For more information on the glossary and acronyms, refer 
% to the link under Glossary and Acronyms in README.md.
\newglossaryentry{LaTeX}{
	name={\LaTeX},
	description={A document preparation system used for
		high-quality typesetting. It is most often used
		for medium-to-large technical or scientific documents.
		\LaTeX is not a word processor!
	}
}

\newglossaryentry{Unix}{
	name={Unix},
	description={A multi-user multi-tasking operating system
		developed by Ken Thompson and Dennis Ritchie in the 
		1970s, which has inspired many features behind Mac OSX 
		and Linux
}
}

\addbibresource{./example.bib}

%%%%%%%%%%%
%  BACK MATTER  %
%%%%%%%%%%%

% The \begin{backmatter} command will create the backmatter environment. This 
% environment includes the formatting required for the body, as well as the the 	
% appendix environment.
\begin{backmatter}

% The appendix environment creates an appendix that will be properly added to 
% table of contents. All figures and tables will be properly formatted to include the 
% appendix letter.
\begin{appendix}{Here's An Appendix}
		
	\lipsum[1-2]
	
\end{appendix}

\begin{appendix}{Here's Another Appendix}
	
	\lipsum[1-2]
		
\end{appendix}

\begin{appendix}{A Figure in An Appendix}
	
	\lipsum[1-2]
	
	% This figure shows proper appendix formatting.
	\begin{figure}[!ht]
		\centering
		\includegraphics[width=5cm]{./stock-image.jpg}
		\caption{It's back!}
	\end{figure}

\end{appendix}

\begin{appendix}{A Table in An Appendix}
        
        \lipsum[1-2]

	% This table shows proper appendix formatting.
        \begin{table}
		\centering
		
		\begin{tabular}{|c|c|c|} \hline
			\textbf{Data} & \textbf{Integer} & \textbf{Float} \\ \hline
			Foo & 1 & 1.1 \\ \hline
		\end{tabular}
		
		\caption{A Table in an Appendix, displaying the correct numbering}
        \end{table}

\end{appendix}

\end{backmatter}

\end{document}
